\selectlanguage{ngerman}
\renewcommand{\abstractname}{Zusammenfassung}
\begin{abstract}
Der Mangel an medizinischen Trainingsdaten erschwert das Training von KI-Modellen in der Radiologie. Synthetische Daten bieten hier einen L�sungsansatz. Diese Arbeit untersucht, ob sich mit einem Diffusionsmodell realistische synthetische 3D-Hirn-MRTs erzeugen lassen. Auf Basis des BraTS-2023-Datensatzes wird ein unconditional Basismodell aufgebaut, in Ablationsstudien analysiert und mit vier Sampling-Verfahren verglichen. Dar�ber hinaus wird eine zweistufige Klasse-zu-Maske-zu-MRT-Pipeline entwickelt, die vollst�ndig synthetische Volumina mit Tumor erzeugt. Die Auswertung erfolgt quantitativ (FID, SSIM, PSNR, LPIPS) und qualitativ gegen den realen Datensatz. Die funktionierenden Modelle erreichen FID-Werte um 25 und erzeugen visuell plausible, anatomisch hirn�hnliche Volumina. Die Pipeline erreicht dabei das Niveau der unconditional Baseline. Als zentrale Einschr�nkung bleibt eine systematische Untersch�tzung der Tumorgr��e um etwa Faktor zwei. Insgesamt zeigt die Arbeit, dass diffusionsbasierte 3D-Synthese ein realistischer Weg ist, um den Datenmangel zu lindern, sofern solche Verzerrungen adressiert werden.
\end{abstract}